\documentclass{report}
\usepackage{amsmath,amssymb}
\setlength{\parindent}{0mm}
\setlength{\parskip}{1em}
\begin{document}
\begin{center}
\rule{6in}{1pt} \
{\large Erin Molloy \\
{\bf Is the ideal approximation operator always ``ideal'' for a particular C/F splitting?}}

University of Illinois at Urbana-Champaign \\ Department of Computer Science
\\
{\tt emolloy2@illinois.edu}\\
Luke Olson\\
Jacob Schroder\end{center}

Given a coarse grid, the ideal prolongation operator is defined by
$\mathbf{P}_\star = \begin{bmatrix} \mathbf{W} & \mathbf{I} \end{bmatrix}^T$,
where the weight matrix, $\mathbf{W} = \mathbf{A}_{FF}^{-1} \mathbf{A}_{FC}$ ,
interpolates a set of fine grid variable ($F$-points) from a set of coarse grid variable
($C$-points), and the identity matrix, $\mathbf{I}$, represents the
injection of $C$-points to and from the coarse grid (Falgout and
Vassilevski, 2004).
In this talk, we consider $\mathbf{P}_\star$, constructed from
both traditional $C/F$ splittings and $C/F$ splittings corresponding to
aggregates, for several challenging problems. We demonstrate the effects of the
$C/F$ splitting on the convergence and complexity of $\mathbf{P}_\star$.
Finally, we argue that $\mathbf{P}_\star$ may be misleading in demonstrating the
``ideal'' nature of interpolation of a given $C/F$ splitting by providing
numerical evidence that hierarchies built using $\mathbf{P}_\star$
converge more slowly
than hierarchies built from alternative prolongation operators with the
same $C/F$ splitting. This is important as we wish to minimize the number
of levels in
a multigrid hierarchy by using a small set of C points for which $\mathbf{P}_\star$
may have poor convergence.


\end{document}
